In open-loop urban navigation without continuous external localization corrections, odometry drift accumulates linearly with cumulative travel distance $s$, progressively corrupting the vehicle's computed line-of-sight bearing toward a target at fixed horizontal distance $D=40$ m. Problem 3 requires determining the maximum travel distance $s_{\max}$ that preserves at least 95\% probability of keeping the wrapped horizontal azimuth pointing error within $\pm\Delta_{\max}=\pm 5^\circ$ (Eq.~\eqref{eq:p3-prob}), and characterizing the feasible sensor-specification region in $(a,b)$ space at $s=80$ m \cite{jcgm2008}.

\subsection{Wrapped Error Propagation Model and Strict Monotonicity Theorem}
By rotational invariance of the isotropic bivariate Gaussian ground-position error $(\delta x,\delta y)\sim\mathcal{N}(0,\sigma(s)^2 I_2)$ with $\sigma(s)=a+bs$ ($a=0.3$ m, $b=0.02$ m/m), we align the $+x$ axis with the true vehicle-to-target vector $(D,0)$. Writing $(\delta x,\delta y)=\sigma(s)(-Z_x,-Z_y)$ with $Z_x,Z_y\stackrel{\mathrm{i.i.d.}}{\sim}\mathcal{N}(0,1)$, the displaced vehicle-to-target vector becomes $(D+\sigma Z_x,\,\sigma Z_y)$, inducing the exact geometric bearing error $\Theta=\operatorname{atan2}(\sigma Z_y,\,D+\sigma Z_x)\in(-\pi,\pi]$. Adding the independent zero-mean Gaussian gimbal zero-offset ($\sigma_{\mathrm{off}}=0.8^\circ$) and control ($\sigma_{\mathrm{ctrl}}=0.6^\circ$) errors yields a combined gimbal perturbation $\epsilon_g\sim\mathcal{N}(0,\sigma_g^2)$ with $\sigma_g=\sqrt{0.8^2+0.6^2}=1.0^\circ$, and total wrapped horizontal pointing error
\[
\Delta\psi = \operatorname{wrap}_{[-\pi,\pi)}\!\bigl(\Theta + \epsilon_g\bigr).
\]
Before solving numerically for $s_{\max}$, one must establish mathematically that the coverage function $P(\sigma(s))=\mathbb{P}(|\Delta\psi|\le\Delta_{\max})$ cannot oscillate or cross $0.95$ multiple times—otherwise a local root-finder could return a spurious threshold. We resolve this rigorously in Theorem~\ref{thm:p3-monotone}.

\begin{keybox}
\begin{theorem}[Strict Monotonicity, Boundary Closure, and Unique Admissible Distance]
\label{thm:p3-monotone}
Let the ground position error be isotropic Gaussian $(\delta x, \delta y) \sim \mathcal{N}(0, \sigma^2 I_2)$ with non-negative spread $\sigma(s) = a + bs$ ($a \ge 0,\; b \ge 0$) at fixed target range $D > 0$, and let the independent gimbal pointing error be $\epsilon_g \sim \mathcal{N}(0, \sigma_g^2)$ with $\sigma_g = \sqrt{\sigma_{\mathrm{off}}^2 + \sigma_{\mathrm{ctrl}}^2} > 0$. Under any forward-wedge tolerance $0 < \Delta_{\max} < \pi/2$ satisfying $P(0) \ge 0.95 > \Delta_{\max}/\pi$, the exact wrapped coverage probability $P(\sigma) = \mathbb{P}\!\left(|\operatorname{wrap}_{[-\pi,\pi)}(\Theta+\epsilon_g)| \le \Delta_{\max}\right)$ extends continuously to $\sigma=0$ with $P(0) = \mathbb{P}(|\operatorname{wrap}_{[-\pi,\pi)}(\epsilon_g)| \le \Delta_{\max})$ and is strictly decreasing in $\sigma$ on $[0, \infty)$, yielding a unique critical spread $\sigma^* \ge 0$ solving $P(\sigma^*) = 0.95$. Consequently:
\begin{enumerate}[leftmargin=1.5em,itemsep=1pt,topsep=2pt,parsep=0pt]
  \item For $b > 0$ and $a \le \sigma^*$ (equivalently $P(a) \ge 0.95$), there exists a unique maximum admissible travel distance $s_{\max} = (\sigma^* - a)/b \ge 0$ solving $P(a + b s_{\max}) = 0.95$.
  \item On the zero-drift boundary $b = 0$, $P(a + 0\cdot s) \equiv P(a)$ is constant in $s$, so coverage stays $\ge 0.95$ for all $s \ge 0$ if and only if $a \le \sigma^*$.
  \item At any fixed travel distance $s_0 > 0$, the exact admissible non-negative sensor domain achieving $P(\sigma(s_0)) \ge 0.95$ is the closed triangular simplex $\{(a,b) \in \mathbb{R}_{\ge 0}^2 : a + s_0 b \le \sigma^*\}$.
\end{enumerate}
\end{theorem}
\end{keybox}

\begin{proof}
For $\sigma > 0$, define the dimensionless signal-to-noise ratio $\kappa = D/\sigma > 0$, so that $\Theta = \operatorname{atan2}(Z_y,\, \kappa + Z_x)$ has distribution depending on $(D,\sigma)$ solely through $\kappa$. Let $\varphi(u)=(2\pi)^{-1/2}e^{-u^2/2}$ and $\Phi(u)=\int_{-\infty}^u\varphi(t)\,\mathrm{d}t$ denote the standard normal density and distribution function.

\emph{Step 1 (Bearing distribution).}
For $0<\alpha<\pi$, the event $|\Theta|\le\alpha$ is $\kappa+Z_x\ge |Z_y|\cot\alpha$, up to a null set. Thus
\[
F_\Theta(\alpha;\kappa)=\int_{\mathbb R}\Phi(\kappa-|y|\cot\alpha)\varphi(y)\,\mathrm{d}y.
\]
The derivative of the integrand in $\kappa$ is positive and bounded by $\varphi(y)/\sqrt{2\pi}$ for every $\kappa$. Dominated convergence therefore gives $0<\partial_\kappa F_\Theta\le 1/\sqrt{2\pi}$. Completing the square gives $\partial_\kappa F_\Theta=\sqrt{2/\pi}\sin\alpha\,e^{-\kappa^2\sin^2\alpha/2}\Phi(\kappa\cos\alpha)$.

\emph{Step 2 (Wrapped acceptance).}
Integrating the periodized Gaussian against $e^{-in\theta}$ unfolds the integral onto $\mathbb R$ and gives Fourier coefficient $q^{n^2}/(2\pi)$, $q=e^{-\sigma_g^2/2}\in(0,1)$ \cite{dlmf}. The Jacobi product gives
\[
g_W(\theta)=\frac1{2\pi}\prod_{n\ge1}(1-q^{2n})
   (1+2q^{2n-1}\cos\theta+q^{4n-2}).
\]
The positive factors yield a positive product since the deviations from one are absolutely summable. Their log derivatives are $-2q^{2n-1}\sin\theta/(1+2q^{2n-1}\cos\theta+q^{4n-2})$, bounded in absolute value by $2q^{2n-1}/(1-q)^2$. The uniform geometric majorant allows termwise differentiation of the log, giving $g_W'<0$ on $(0,\pi)$. For the even acceptance kernel $I(\theta)=\int_{\theta-\Delta_{\max}}^{\theta+\Delta_{\max}} g_W(u)\,\mathrm{d}u$, the continuous integrand gives $I'=g_W(\theta+\Delta_{\max})-g_W(\theta-\Delta_{\max})<0$ there, since the first argument is farther from zero on the circle.

\emph{Step 3 (Continuity at $\sigma=0$, chain rule, and closed-form boundary simplex).} Conditioning on $|\Theta|$ and integrating by parts on $[0,\pi]$ (using $F_\Theta(0;\kappa)=0$ and $F_\Theta(\pi;\kappa)=1$) yields
\[
P(\sigma) = \int_0^\pi I(\theta)\,\mathrm{d}F_\Theta(\theta;\kappa)
= I(\pi) + \int_0^\pi \bigl(-I'(\theta)\bigr)\,F_\Theta(\theta;\kappa)\,\mathrm{d}\theta.
\]
For $\sigma\downarrow0$, $F_\Theta(\theta;D/\sigma)\uparrow1$ at interior angles, so $P(\sigma)\to I(0)$ by dominated convergence. For $\sigma\uparrow\infty$, the isotropic bearing tends to uniform and coverage tends to $\Delta_{\max}/\pi$. On $[\sigma_1,\sigma_2]\subset(0,\infty)$, the derivative of the integral's integrand in $\sigma$ is strictly negative at interior angles and its absolute value is bounded by $D|I'(\theta)|/(\sigma_1^2\sqrt{2\pi})$. This envelope is integrable since $I'$ is continuous on $[0,\pi]$. Differentiating shows $P'(\sigma)<0$, so the 95\% level is crossed once under the stated bounds, and $P(a+bs)\ge0.95$ is equivalent to $a+bs\le\sigma^*$.
\end{proof}

\subsection{Analytical-Statistical Dual Closure and Sensor Specification Domain}
With global uniqueness established by Theorem~\ref{thm:p3-monotone}, we compute $s_{\max}$ via independent conditional Simpson quadrature (192 nodes truncated at $\pm 10\sigma$, tail mass $<10^{-22}$, cross-checked at 96 and 384 nodes), yielding $s_{\max}=\BIntegratedCrossing{}$ m at critical position spread $\sigma_{\mathrm{loc}}(s_{\max})=\BSigmaCrossing{}$ m, strictly bracketed within $[\BBracketLow{},\,\BBracketHigh{}]$ m (width $2\times 10^{-4}$ m; superseding the legacy $\BTravelThreshold{}$ m baseline). For large distances ($s\ge 240$ m), an explicit union bound certifies $P(\sigma(s))\le \BFarBound{}$, further confirming asymptotic decay far below 0.95.

To close the loop statistically without relying on quadrature alone, we executed three fixed-seed Monte Carlo probes ($n=\BMCBeforeSamples{}$ strictly independent trials per probe, $6\times 10^6$ independent trials total) and constructed simultaneous $\BMCConfidence\%$ Wilson score confidence intervals \cite{wilson1927,brown2001} with Bonferroni multiplicity correction ($\alpha/3$). As listed in Table~\ref{tab:p3} and plotted in Figure~\ref{fig:p3}b and Figure~\ref{fig:p3}d, the lower probe at $\BMCBeforeDistance{}$ m achieves empirical coverage $\BMCBeforeEstimate{}$ with simultaneous interval $[\BMCBeforeLow{},\,\BMCBeforeHigh{}]$ lying strictly above $0.95$, whereas the upper probe at $\BMCAfterDistance{}$ m drops to $\BMCAfterEstimate{}$ with interval $[\BMCAfterLow{},\,\BMCAfterHigh{}]$ lying strictly below $0.95$. At the benchmark check distance of $s=\BMCLimitDistance{}$ m, coverage degrades to $\BMCLimitEstimate{}$ ($[\BMCLimitLow{},\,\BMCLimitHigh{}]$), confirming that nominal hardware ($a=0.3$ m, $b=0.02$ m/m) fails the 95\% requirement at 80 m.

Variance ablation clarifies the dominant physical driver: at $s=0$ m, position error accounts for only $15.59\%$ of total angular variance (zero-offset $54.02\%$, control $30.39\%$), whereas by $s_{\max}=\BIntegratedCrossing{}$ m, self-localization drift accounts for $84.61\%$ of angular variance (zero-offset $9.85\%$, control $5.54\%$; shifting to $88.10\%$, $7.61\%$, and $4.28\%$ respectively at 80 m), and removing position error alone ($\sigma=0$) restores coverage to $P(0)=0.9999994$. By Theorem~\ref{thm:p3-monotone}(3), meeting the 95\% specification at $s_0=80$ m is equivalent to the exact closed-form non-negative parameter simplex
\[
a \ge 0,\qquad b \ge 0,\qquad a + 80\,b \le \sigma^* = \BSigmaCrossing{}\text{ m},
\]
with vertices $(0,0)$, $(\BSigmaCrossing{}\text{ m},\,0)$, and $(0,\,0.020459\text{ m/m})$ (annotated directly in Figure~\ref{fig:p3}c)—defining the complete admissible odometry procurement envelope without numerical contour approximation.
